\documentclass[11pt,a4paper]{article}
\usepackage[margin=26mm]{geometry}
\usepackage{amsmath,amssymb,amsthm,booktabs}
\usepackage[hidelinks]{hyperref}
\usepackage{microtype}
\newcommand{\cut}{c}
\newcommand{\RR}{\mathbb R}
\newcommand{\eff}{\operatorname{Eff}}
\newtheorem*{theorem}{Theorem}
\setlength{\parindent}{0pt}
\setlength{\parskip}{5pt plus 1pt minus 1pt}
\setlength{\abovedisplayskip}{6pt}
\setlength{\belowdisplayskip}{6pt}
\hypersetup{pdftitle={The three least-order counterexamples to a Cheeger bound for Pareto-efficient vectors}}
\begin{document}
\begin{center}
{\Large\bfseries The three least-order counterexamples to a Cheeger bound\\[3pt]
for Pareto-efficient vectors\par}
\smallskip
{\small Research note \(\cdot\) 25 November 2026}
\end{center}

\textbf{Abstract.}
For a connected simple graph with ten vertices and eighteen edges, the edge
expansion is \(4/3\) and the minimum total weight of a Pareto-efficient point
of the small-cut polytope is \(6\). This disproves the floor-corrected lower
bound proposed by jinni. Every graph on two through nine vertices satisfies
the bound for every efficient real vector. An exact census shows that
precisely three isomorphism classes violate it at order ten. Their edge
counts are 18, 18, and 21.

\section*{The question and the graph}
For a finite simple graph \(G=(V,E)\), write \(n=|V|\ge2\) and
\(k=\lfloor n/2\rfloor\). Let \(\cut(S)\) count the edges with exactly one endpoint
in \(S\), and write \(x(S)=\sum_{v\in S}x_v\). Define
\[
P_G=\{x\in\RR_{\ge0}^{V}:x(S)\le\cut(S)\text{ whenever }|S|\le k\},
\qquad h(G)=\min_{1\le |S|\le k}\frac{\cut(S)}{|S|}.
\]
A vector \(x\in P_G\) is Pareto-efficient if every \(y\in P_G\) with
\(y\ge x\) coordinatewise equals \(x\). Write \(\rho(G)=\min_{x\in\eff(P_G)}x(V)\).
The revised question asks whether \(\rho(G)\ge k h(G)\) [1,2].
DSM's examples for the original coefficient \(|V|/2\) have odd order [3].

\begin{theorem}
The graph whose graph6 string is \textnormal{\texttt{IJOG[BHvO}} has
\(\rho(G)=6<20/3=\lfloor |V|/2\rfloor h(G)\).
Its order is the least possible. Among graphs of that order there are two
counterexamples with eighteen edges and one with twenty-one edges.
\end{theorem}

Take three disjoint triangles \(Q_i\). Join the three distinguished vertices
in the middle column of the table to form a fourth triangle \(D\).
Join vertex \(9\) to the indicated pair in each row.
\begin{center}
\begin{tabular}{ccc}\toprule
Vertices of \(Q_i\) & Vertex of \(D\) & Neighbors of \(9\) in \(Q_i\)\\ \midrule
\(\{0,7,8\}\) & \(8\) & \(0,7\)\\
\(\{1,2,3\}\) & \(1\) & \(1,3\)\\
\(\{4,5,6\}\) & \(4\) & \(4,5\)\\ \bottomrule
\end{tabular}
\end{center}
These four edge-disjoint triangles and six additional edges give eighteen edges.
Every vertex is within distance two of \(9\).

\section*{The cut bounds}
Fix \(S\) with \(|S|\le5\). Write \(a_i=|S\cap Q_i|\), \(t=|S\cap D|\), and
let \(m_i\) count the neighbors of \(9\) in \(S\cap Q_i\). Each \(Q_i\)
contributes \(a_i(3-a_i)\) crossing edges, and \(D\) contributes \(t(3-t)\).
If \(9\notin S\), the remaining contribution is \(\sum_i m_i\).
If \(9\in S\), it is \(6-\sum_i m_i\).

\textit{Sets omitting \(9\).}
If \(a_i\le2\), then \(a_i(3-a_i)+m_i\ge 3a_i/2\).
Indeed one selected vertex has two internal edges leaving it, and two selected
vertices include at least one neighbor of \(9\). Thus if \(S\) contains no full
triangle \(Q_i\), its boundary is at least \(3|S|/2\).

\newpage
If \(S\) contains a whole \(Q_i\), put \(j=|S|-3\le2\). This \(Q_i\)
contributes two spokes and contains one vertex of \(D\). For \(t=1,2\),
the triangle \(D\) contributes two more edges and the remaining rows contribute
at least \(3j/2\). Hence \(\cut(S)\ge4+3j/2\ge4|S|/3\).
For \(t=3\), the extra vertices must be the two other distinguished vertices.
They each contribute two edges of their own row, and at least one is a neighbor
of \(9\). Thus \(\cut(S)\ge2+4+1=7\) for this five-set.
Altogether \(\cut(S)\ge4|S|/3\) for \(|S|\le5\) omitting \(9\).
In particular every such five-set has at least seven crossing edges.

\textit{Sets containing \(9\).}
For a five-set, take its complement and apply the preceding seven-edge bound.
For \(|S|\le4\), at most three vertices besides \(9\) are selected. If they fill
a row, the four spokes to the other rows and the two crossing edges of \(D\)
give \(\cut(S)=6\). Otherwise \(a_i\le2\), so
\(a_i(3-a_i)-m_i\ge0\) in each row. The cut formula gives \(\cut(S)\ge6\).
In both cases \(\cut(S)\ge4|S|/3\). Each individual row \(Q_i\) has three
vertices and boundary four, so \(h(G)=4/3\).

\section*{The minimum over efficient real vectors}
Set \(x_9=6\) and \(x_v=0\) elsewhere. Every admissible set containing \(9\)
has boundary at least six, which proves feasibility. Each \(U_i=Q_i\cup\{9\}\)
has size four and boundary six. If a feasible real vector \(y\) dominates \(x\),
then
\[
0\le\sum_{v\in U_i}(y_v-x_v)\le \cut(U_i)-x(U_i)=0 .
\]
All summands are nonnegative, so \(y=x\) on each \(U_i\). The three sets cover
\(V\), proving efficiency.

Conversely, if a feasible \(z\) has \(z(V)<6\), add \(6-z(V)>0\) to its
coordinate at \(9\). Its new sum on any set containing \(9\) is at most six
by nonnegativity, and every other sum is unchanged. This produces a distinct
feasible vector dominating \(z\). Thus every efficient vector has mass at
least six, and \(x\) attains this bound. The violation is \(20/3-6=2/3\).

\section*{Why a smaller counterexample would need a special cover}
Suppose that an efficient real vector \(z\) has total \(w<kh(G)\), so
\(h(G)>0\). Call an admissible set \(S\) \emph{tight} when \(z(S)=\cut(S)\).
Every vertex belongs to a tight set. Otherwise its coordinate could be
increased by the minimum positive slack among the finitely many
constraints containing it. The inclusion-maximal tight sets therefore
cover \(V\). Each has size strictly less than \(k\), since
\(h(G)|S|\le\cut(S)=z(S)\le w\).

Take two distinct maximal tight sets \(A,B\) and put \(U=A\cup B\),
\(I=A\cap B\). The cut function is submodular since the difference
\(\cut(A)+\cut(B)-\cut(U)-\cut(I)\) counts twice the edges between
\(A\setminus B\) and \(B\setminus A\). Feasibility on \(I\) gives
\[
  \cut(U)\le\cut(A)+\cut(B)-\cut(I)
       \le z(A)+z(B)-z(I)=z(U)\le w.
\]

\newpage
\noindent
If \(|U|\le k\), feasibility makes \(U\) tight, contradicting maximality.
If \(|U|=n-k\), the complement has size \(k\), so
\(\cut(U)=\cut(V\setminus U)\ge kh(G)>w\). The sets must intersect,
since disjoint sets would have combined mass at least
\(h(G)(|A|+|B|)>kh(G)\). Thus every distinct pair satisfies
\[
 k<|A\cup B|\le 2(k-1)-1,\qquad |A\cup B|\ne n-k.
\]
There are at least two maximal tight sets since each has size less than \(k\).
For \(2\le n\le9\), the displayed conditions force \(k=4\) and
\(|A\cup B|=5\). They exclude \(n=9\) since \(9-4=5\). At \(n=8\)
every maximal tight set must be a triple, and every two triples must meet
in exactly one vertex.

\section*{Eight vertices cannot be covered by such triples}
Three distinct triples exist because two cover at most five vertices.
Choose three of them.
The three together cover at most seven, since the second and third
each meet the first. Take a fourth triple through an uncovered point.
Its two remaining positions must meet all three earlier triples.
One point \(p\) therefore belongs to three distinct triples.

These three triples through \(p\) cover exactly seven points because their
three pairs outside \(p\) are disjoint. Let \(W\) cover the eighth point.
If \(p\notin W\), meeting all three triples requires one vertex from
each of the three disjoint pairs, leaving no position for the eighth point.
If \(p\in W\), both its other vertices must be outside the seven to avoid
a second intersection. There is only one such vertex. This contradiction
proves the lower bound for all real vectors. Together with the graph
above it proves that the minimum order is ten.

\section*{The three order-ten classes}
In McKay's labels [4], all three graphs have the triangles
\(\{0,3,6\},\{1,4,7\},\{2,5,8\}\) and \(\{0,7,8\}\).
Their only other edges join \(9\) to the neighbors listed below.
\begin{center}
\begin{tabular}{l l r c c}\toprule
graph6 & Neighbors of \(9\) & Edges & \(h\) & \(\rho\)\\ \midrule
\verb!ICOceQd{W! & \(0,1,2,3,7,8\) & 18 & \(4/3\) & 6\\
\verb!ICOceQd|O! & \(0,1,2,3,5,7\) & 18 & \(4/3\) & 6\\
\verb!ICOceQd~w! & \(0,1,\ldots,8\) & 21 & \(5/3\) & 8\\ \bottomrule
\end{tabular}
\end{center}
In each case \(x=\rho e_9\) is feasible and the first three triangles,
each augmented by \(9\), are tight and cover \(V\). Every admissible
cut through \(9\) has capacity at least \(\rho\), so the deficit argument
above proves the exact minimum. The finite inequalities and attained
expansions are checked in Lean. Only the second 18-edge graph has a
degree-four vertex. The map with successive values
\((2,7,4,1,0,3,6,5,8,9)\) carries the original example to this second graph.

\newpage
\section*{An exact certificate of exhaustiveness}
In a putative violation, each maximal tight set \(A\) has size 3 or 4.
Indeed there is another such set \(B\), with \(|B|\le4\), intersecting
\(A\) and satisfying \(|A\cup B|>5\). Also \(\cut(A)<5h\). Every pair
of distinct maximal tight sets must satisfy the sharper inequality
\[
 \cut(A)+\cut(B)-\cut(A\cap B)
 \le z(A)+z(B)-z(A\cap B)=z(A\cup B)\le w<5h .
\]
For each graph the certificate specifies an attaining cut \(T\).
All comparisons with \(5h=5\cut(T)/|T|\) use integer cross-products.
Start with the three- and four-sets of capacity below this threshold.
A negative proof tree names an uncovered vertex at each node. For
every available set through that vertex there is a child, with its
vertices removed and the pool restricted to compatible sets meeting
the remainder. A pivot with no available set closes a branch. A
putative tight cover would follow children to an empty remainder,
where no pivot can be supplied. Each supplied tree is checked in full.

Of the 11,716,571 connected representatives, the initial pool misses
a vertex for 9,885,683. The pair constraints exclude a further
1,830,885. The three remaining graphs are exactly the realizations
in the table. Among the 774,628 graphs with fewer than 18 edges every
one was excluded. Of the 561,106 graphs with 18 edges precisely two
survive. A disconnected graph of order ten has a nonempty component
with at most five vertices, so its Cheeger constant is zero.

Completeness of the input was checked independently. Ordered partition
refinement and individualization gave distinct canonical adjacency
strings for every input. The number of labelings realizing the chosen
canonical string gives the automorphism order. Summing
\(10!/|\operatorname{Aut}(G)|\) gave \(34{,}496{,}488{,}594{,}816\).
This equals the number \(C_{10}\) of connected labeled graphs. Partitioning
all labeled graphs according to the component of vertex 0 gives
\[
C_n=2^{\binom n2}-
  \sum_{j=1}^{n-1}\binom{n-1}{j-1} C_j\,2^{\binom{n-j}{2}},
\qquad C_1=1.
\]
This establishes the classification up to isomorphism. \hfill\(\square\)

{\small\textbf{Verification and provenance.}
The original graph and its witness were supplied in the review archive.
Lean checks the real-vector exclusion rule, the witnesses and the
nonisomorphism result. Two C++ programs produce and replay the order-ten
census with different cut calculations. The project includes all input
graphs, refutation trees, automorphism orders and a Python checker [5].}

{\small
\textbf{Sources}\par
[1] jinni, \emph{Conjecture involving the Cheeger constant of a graph},
\href{https://mathoverflow.net/questions/506070}{MathOverflow 506070},
18 December 2025. Conjecture B was added in
\href{https://mathoverflow.net/revisions/c774f5bd-ab71-4085-8a45-7df5d7b6ddec/view-source}{revision 5},
20 December 2025.\par
[2] jinni, \emph{Conjecture involving the Cheeger constant of a graph},
\href{https://math.stackexchange.com/questions/5118288}{Mathematics Stack Exchange 5118288},
5 January 2026.\par
[3] DSM, \href{https://mathoverflow.net/a/506124}{answer to MathOverflow 506070},
19 December 2025.\par
[4] Brendan D. McKay, \href{https://users.cecs.anu.edu.au/~bdm/data/graphs.html}{\emph{Combinatorial Data: Simple graphs}},
connected graph6 catalogs, orders 2--10, accessed 25 November 2026.\par
[5] \href{https://github.com/Kazeshima/solved-problems/tree/main/problems/cheeger-pareto}{\emph{Cheeger--Pareto proofs, census certificates and replay programs}}, project repository.}
\end{document}
